\chapter{The shape of neighbourhood vulnerability}\label{chapter:vulnerability}

\section{Introduction}

The map resolved heat--mortality risk to the neighbourhood and showed that it varies there. It did not say what a neighbourhood is made of that makes it vary. The risk at a dissemination area was read from a vulnerability profile, but the profile entered the prediction as a single composite quantity, and a composite hides its parts. We know the risk differs from one neighbourhood to the next; we do not yet know along how many distinct lines it differs, or what those lines are.

This chapter reads the profile. We build the vulnerability indicators from what a neighbourhood measurably is, find how many distinct kinds of vulnerable neighbourhood the country holds, and say what each kind is made of. The number of kinds is not the number of indicators. A poorer neighbourhood tends also to be a less green one, and a renting one; indicators that rise and fall together across the country are not separate kinds of vulnerability but several readings of one. The work of the chapter is to find how many kinds remain once that redundancy is taken out.

\section{What the profile is}

The profile is built from what a neighbourhood is, never from what happens to it. Its inputs are measured at the dissemination area: the age of its residents, the condition of its housing, the heat of its surfaces. None of these is a death count. A profile built from the outcome would explain that outcome by restating it; this one cannot, because the outcome is nowhere in it.

This constraint is also what makes the profile useful. The previous chapter could read a curve for every dissemination area only because the profile was defined for every dissemination area, including the ones whose death counts can support no fit of their own. A profile that needed local mortality to compute would fail at exactly the neighbourhoods the design exists to reach. Building it from inputs alone is therefore the honest choice and the one that generalizes; the same property that keeps the profile from borrowing against the outcome is what lets it be read everywhere the inputs exist.

The raw material is a set of indicators measured at the dissemination area. They fall into a few families: the physical environment a neighbourhood sits in, the climate its built surfaces generate, the makeup of the people who live there, and the conditions they live in. Table~\ref{tab:indicators} lists the candidate set. Every one is drawn from public census and environmental data that already exists at the dissemination area, which is what lets the profile resolve to the neighbourhood at all. The final set is fixed once the variable request returns; what matters here is that the indicators are measured at the scale the profile is read.

\begin{table}[ht]
\centering
% [Candidate indicator set; final selection fixed once the variable
%  request returns. Columns: indicator, domain, pathway to heat risk.]
\caption{Candidate vulnerability indicators, by domain and pathway to heat-mortality risk.}
\label{tab:indicators}
\end{table}

One choice in how the indicators are combined matters enough to state here. They are measured on different scales: a temperature in degrees, a count in people, a proportion between zero and one. We standardize each to the same scale before combining them, so that a component reflects how indicators move together and not how large their raw numbers happen to be. In the language of the method, the analysis runs on the correlation matrix of the indicators rather than their covariance. A variable measured in larger units would otherwise dominate the result for no reason but its units.

\section{The dimensions}

Principal component analysis takes the standardized indicators and finds the directions along which neighbourhoods differ most. Each direction is a weighted blend of the original indicators, and the neighbourhoods spread out along it more than along any other. The first direction is the one that captures the most variation across the country; the second captures the most of what is left, and so on down. What the method returns is a set of these directions, ordered by how much of the variation each one carries.

Behind this is not a sequence of steps but a single operation. The correlation matrix of the indicators is decomposed into its eigenvectors and eigenvalues,
%
\begin{equation}\label{eq_eigen}
R = V \Lambda V^{\top},
\end{equation}
%
where $R$ is the correlation matrix of the indicators, the columns of $V$ are the eigenvectors, and $\Lambda$ is the diagonal matrix of their eigenvalues. Each eigenvector is one of the directions, and its eigenvalue is the variation the neighbourhoods carry along it. The directions come out ordered because the eigenvalues do, the largest first. Given $R$, this decomposition is unique, with nothing estimated and nothing drawn at random. Nothing in it comes from the death counts: the directions are fixed by how the indicators vary together, and the indicators were measured before a single death entered the analysis.

The eigenvectors have a geometric meaning. Place each neighbourhood as a point in a space with one axis per indicator, positioned by its values. The neighbourhoods form a cloud in that space, and the cloud has a shape. The directions of greatest variation are the longest axes of that cloud, the ones along which it is most stretched.

Whether those axes mean anything depends on the shape. A cloud stretched far along one axis has a clear longest direction: every other direction is shorter, and there is no ambiguity about which way it runs. A cloud that is round has no longest direction at all. Every axis through it is the same length, so the choice of axis is arbitrary, and a line drawn one way is no more the cloud's own than a line drawn another. Real clouds sit between these, and it is a matter of degree: the further a cloud is from round, the more its longest axis is its own, and not an artifact of where the line was drawn.

This is what makes the shape matter downstream. When one axis is much longer than the next, it is stable: change the indicators a little, drop one or add another, and the longest axis still points essentially where it did. When two axes are nearly the same length, they are not stable: a small change in the indicators can swap their order or rotate them into each other, and a direction that seemed to be the cloud's own turns out to have depended on the exact inputs. The two cases are separated by one quantity, the gap between an eigenvalue and the next. A wide gap is a stable direction; a narrow gap is two directions the data cannot tell apart.

The ordered eigenvalues, taken together, are the spectrum, and the spectrum is what the rest of the chapter reads. Its values give the length of each direction; the gaps between them give which directions are the cloud's own and which the data cannot resolve. How many directions stand clear of the rest is the question of how many distinct kinds of vulnerable neighbourhood the country holds, and it is the question §4 takes up next.

\section{How many kinds}

The question this section answers is how many distinct kinds of vulnerable neighbourhood the country holds. It is not how many indicators there are. The indicators were never meant to be counted; many of them measure the same underlying thing from different angles, and counting them would count that thing several times over. The number we are after is the number that remains when indicators that move together are recognized as one: the count of genuinely separate patterns in how Canadian neighbourhoods are vulnerable. We call it $k$.

The spectrum is where $k$ is read. A kind of vulnerability is a direction that stands clear of the rest, and in the spectrum that means an eigenvalue separated from the next by a wide gap. Count the directions that stand clear this way and you have counted the kinds. The count stops where the gaps close: past that point the eigenvalues run close together, and a near-equal eigenvalue is indicators blurring into a direction the data cannot separate.

We can say in advance what we expect the spectrum to show. The literature on heat vulnerability returns to three recurring themes: material deprivation, age, and social isolation.% [CITATION DEBT — the recurring-themes claim needs 5--10 references establishing deprivation/age/isolation as the field's standing heat-vulnerability axes; this is earned in the Introduction (written last), which must do the literature work that licenses stating it flat here. Until then this sentence is unbacked. Klinenberg among intended sources.]
If these are the axes along which neighbourhoods differ, the spectrum should separate into three clear directions before the gaps close. The spectrum either confirms this, by separating where the themes say it should, or breaks it, by separating into more directions or fewer. Which it does is the first thing the spectrum settles that the themes by themselves only suggest.

Computed on the real indicators, the spectrum separates into $k$ directions that stand clear of the rest, and then a tail of eigenvalues that run too close together to separate. The value of $k$, the eigenvalues, and the gaps between them enter here once the indicators are computed on the assembled data; the present draft carries the machinery that reads them, not the values themselves. Figure~\ref{fig:spectrum} shows the spectrum, with the gaps marked where the directions stop standing clear.

\begin{figure}[ht]
\centering
% [Spectrum plot enters here once the indicator set is assembled and
%  the correlation matrix is decomposed: ordered eigenvalues on the
%  vertical axis, component index on the horizontal, gaps between
%  consecutive eigenvalues marked; the wall where separation collapses
%  drawn explicitly.]
\caption{The eigenvalue spectrum of the vulnerability indicators. The gaps between consecutive eigenvalues are marked; $k$ is the number of directions that stand clear of the rest.}
\label{fig:spectrum}
\end{figure}

Each direction that stands clear can be named, and we name it after the indicators that define it. A direction is a weighted blend of all the indicators, but a few carry most of the weight, and those few are what the direction is. The naming follows from them and adds nothing of its own: the first dimension loads most heavily on [its dominant indicators], and we refer to it as the [name] dimension, after those loadings; the second loads on [its dominant indicators], and we name it the [name] dimension on the same basis. The dominant loadings for each clear direction enter here from the loading table once the indicators are assembled. Table~\ref{tab:loadings} reports them.

\begin{table}[ht]
\centering
% [Loading table enters here once the correlation matrix is decomposed:
%  rows are the indicators, columns are the retained directions, entries
%  are the loadings (eigenvector weights); dominant loadings per direction
%  highlighted. The tail directions are shown with no dominant loadings,
%  which is why they are not named.]
\caption{Indicator loadings on the retained directions. Each direction is named after the indicators that load most heavily on it; the tail directions carry no dominant loadings and are left unnamed.}
\label{tab:loadings}
\end{table}

The prior was three. The spectrum either met it or did not, and the comparison is the finding. [Whether the country resolved into the three kinds the literature anticipated, fewer if the themes turned out to bundle into one, or more if they came apart, enters here once the spectrum is computed; the value of $k$ and its distance from the prior are what the finding reports.] The themes name the axes a neighbourhood might vary along; the spectrum says how many of them it separates into, and that count is the kinds of vulnerable neighbourhood the country holds.

What the spectrum gives is a description: how many kinds of vulnerable neighbourhood the country holds, and what each is made of. It does not say why the country took that shape rather than another, and this chapter does not ask. The reason a particular set of kinds emerged, and not some other, is a question about the forces that built Canadian neighbourhoods, and it is taken up in the discussion. 

\section{Naming and its limit}

A dimension can be named only as far as its eigenvalue stands clear of its neighbours. Where two eigenvalues are nearly equal, the data has not picked out a direction; it has picked out a plane, within which any direction does as well as any other. There is nothing to name there, because the data has not separated one thing from the next. We name the directions the spectrum sets apart, and we stop where it stops setting them apart.

This limit is not a threshold we chose. Nameability is a matter of degree, set by how wide the gap is, and a wide gap and a narrow one shade into each other with no natural line between them. We could draw a line anywhere and call everything above it nameable, but the line would be ours, not the data's, and an examiner would be right to ask why it fell there and not a little higher or lower. So we draw no line. We report the gaps as the spectrum gives them and let nameability fall off where the data lets it, naming confidently where the gap is wide, declining where it is narrow, and saying plainly that the cases between are between.

There is a standing objection to naming anything from principal components: change the set of indicators that goes in, and the loadings can shift, so the dimensions you name are an artifact of the inputs you happened to choose [CITATION — Conlon, PCA input-set instability — Verify Friday]. The objection is real, and the spectrum is already the answer to it. Where an eigenvalue stands well clear of its neighbours, its direction is stable under exactly this kind of change: drop an indicator or add one, and a well-separated direction still points where it did, which is what the gap on the spectrum shows. Where eigenvalues are near-equal, the direction does shift under that change, but those are the directions we already declined to name.

Two choices keep the naming honest, and both are the data's, not ours. The components are taken from the correlation matrix, so no indicator earns a dimension by the size of its units. The number of components carried into the model is set by minimizing the Akaike information criterion of the meta-regression, so the count is the one the pooled fit supports and not one we preferred. What survives both is a small set of named dimensions, each standing clear on the spectrum and built from indicators that load on it together. They say what the kinds of vulnerable neighbourhood are. They do not say whether those kinds fall the same way for everyone who lives in them.

\section{What this leaves open}

This chapter built the vulnerability profile and read its shape. The profile is assembled from what a neighbourhood is, never from its deaths; it resolves into a countable set of distinct kinds of vulnerable neighbourhood; and each kind that stands clear can be named from the indicators that define it. What the chapter establishes is the shape of vulnerability across the country: how many kinds there are, and what each is made of. What it does not yet do is say who that shape is for.

The shape was read across all ages at once. Every dimension, every named kind, describes a neighbourhood as a whole, averaged over the young and the old who live in it together. But age is among the strongest things that decide who heat kills,% [CITATION—age as a leading modifier of heat mortality. Stated flat here, to be earned in the Introduction (written last), which establishes age as a standing modifier so the body need not cite it. If the Introduction does not carry it, back here — and vary from gasparrini2022smallarea, which Ch3 already leans on for this claim.]
and the previous chapter standardized it away on purpose, so that the variation on the map would be the neighbourhood's and not an accident of which neighbourhoods happen to be older. That was the right choice for reading the neighbourhood signal cleanly. It also means the shape we have read is the shape for everyone together, and not yet the shape for the people heat most threatens.

The next chapter takes age back up. It splits the map by age group and reads the geography of vulnerability separately for each, restoring the distinction the standardization folded away. Where this chapter asked what kinds of vulnerable neighbourhood the country holds, the next asks whether the same neighbourhoods are the dangerous ones for the old as for the young, or whether the map redraws itself depending on whose risk it is showing.